%%%%%%%%%%%%%%%%%%%%%%%%%%%%%%%%%%%%%%%%%%%%%%%%%%%%%%%%%%%%%%%%%%%%%%%
% BAB 3
%%%%%%%%%%%%%%%%%%%%%%%%%%%%%%%%%%%%%%%%%%%%%%%%%%%%%%%%%%%%%%%%%%%%%%%

\mychapter{3}{BAB 3 TINJAUAN PUSTAKA}


\section{Jaringan Wi-Fi Enterprise}

\subsection{Karakteristik Wi-Fi Enterprise}

Wi-Fi \textit{enterprise} merupakan implementasi jaringan WLAN yang dirancang untuk menyediakan akses jaringan nirkabel pada lingkungan dengan jumlah pengguna dan perangkat yang relatif besar serta kebutuhan pengelolaan jaringan yang lebih kompleks. Implementasinya umumnya melibatkan beberapa \textit{access point} (AP) yang ditempatkan pada area yang berbeda untuk memperluas cakupan layanan jaringan nirkabel, seperti pada lingkungan perkantoran, kampus, pusat perbelanjaan, maupun fasilitas publik \parencite{SEQUEIRA_BUILDING_2017}.

Penggunaan beberapa AP dengan area cakupan yang berdekatan juga memungkinkan perangkat pengguna mempertahankan konektivitas ketika berpindah dari satu area ke area lainnya melalui mekanisme \textit{roaming}. Pada WLAN \textit{enterprise}, kemampuan tersebut diperlukan untuk menjaga kontinuitas layanan ketika perangkat berpindah antar-AP dalam jaringan yang sama \parencite{GERACI_WIFI_2025}. Namun, implementasi beberapa AP pada area yang saling berdekatan memerlukan pengelolaan parameter radio, seperti \textit{channel} dan \textit{transmit power}, untuk mengurangi interferensi dan mempertahankan kualitas layanan.

Peningkatan jumlah AP dan kebutuhan untuk mengelola berbagai aspek jaringan tersebut menyebabkan konfigurasi dan pengelolaan WLAN \textit{enterprise} menjadi lebih kompleks. Untuk menangani kompleksitas tersebut, WLAN \textit{enterprise} umumnya menerapkan mekanisme pengelolaan secara terpusat. Berbeda dengan pendekatan terdistribusi yang menjalankan mekanisme kontrol pada masing-masing AP, pendekatan terpusat menggunakan \textit{controller} untuk mengonfigurasi dan mengelola AP berdasarkan kebijakan dan keputusan yang ditentukan secara terpusat \parencite{DEZFOULI_REVIEW_2019}.

Aspek keamanan juga menjadi salah satu karakteristik penting dalam implementasi jaringan Wi-Fi \textit{enterprise}. Salah satu mekanisme yang digunakan adalah IEEE 802.1X yang menyediakan autentikasi akses jaringan menggunakan model \textit{supplicant}, \textit{authenticator}, dan \textit{authentication server}. Dalam implementasinya pada jaringan WLAN, EAP (\textit{Extensible Authentication Protocol}) digunakan sebagai kerangka untuk mendukung berbagai metode autentikasi, sedangkan \textit{access point} atau perangkat jaringan dapat berperan sebagai \textit{authenticator} yang meneruskan proses autentikasi kepada \textit{authentication server} \parencite{RFC3748,RFC3580}. Penggunaan server autentikasi juga memungkinkan pengelolaan autentikasi pengguna secara terpusat dan dapat digunakan untuk menentukan hak akses jaringan, termasuk penempatan perangkat ke VLAN tertentu berdasarkan hasil autentikasi \parencite{RFC3580}.

Selain autentikasi pengguna, keamanan WLAN juga mencakup perlindungan terhadap \textit{management frame}. Pada jaringan IEEE 802.11, beberapa \textit{management frame}, seperti \textit{deauthentication} dan \textit{disassociation}, dapat dipalsukan oleh pihak yang berada dalam jangkauan jaringan. Protected Management Frames (PMF) yang diperkenalkan melalui IEEE 802.11w menyediakan perlindungan kriptografis terhadap sejumlah \textit{management frame} tersebut sehingga dapat mengurangi risiko serangan berbasis pemalsuan \textit{management frame} \parencite{CISCO_80211W_2024}.

Berdasarkan karakteristik tersebut, implementasi Wi-Fi \textit{enterprise} tidak hanya berfokus pada penyediaan akses jaringan nirkabel, tetapi juga mencakup pengelolaan terpusat, konfigurasi infrastruktur, pengaturan parameter radio, hingga konfigurasi keamanan untuk memastikan jaringan dapat berfungsi sesuai dengan kebutuhan.

\subsection{Arsitektur WLAN Enterprise}

Arsitektur WLAN \textit{enterprise} terdiri atas beberapa komponen yang saling terhubung, meliputi \textit{access point} (AP), \textit{wireless LAN controller} (WLC), dan infrastruktur jaringan untuk menyediakan akses nirkabel serta mendukung pengelolaan secara terpusat. Pada arsitektur berbasis \textit{controller}, fungsi kontrol dan penerusan trafik pengguna dapat dipisahkan. RUCKUS SmartZone menerapkan pemisahan tersebut melalui \textit{control and management plane} pada vSZ dan \textit{data plane}, dengan trafik pengguna dapat diteruskan melalui \textit{local breakout} atau \textit{tunneling} menuju \textit{data plane} \parencite{RUCKUS_PRODUCT_GUIDE_2018}.

\section{Komponen Infrastruktur Wi-Fi}

\subsection{Access Point}

Subbab ini membahas fungsi \textit{access point}, prinsip kerja perangkat, serta perannya dalam menyediakan layanan akses jaringan nirkabel kepada pengguna.

\subsection{Wireless LAN Controller (WLC)}

Subbab ini menjelaskan fungsi \textit{wireless LAN controller} sebagai pusat pengelolaan jaringan Wi-Fi, termasuk konfigurasi perangkat, distribusi kebijakan, dan pengawasan jaringan.

\subsection{Access Switch}

Subbab ini membahas peran \textit{access switch} sebagai penghubung perangkat jaringan serta implementasi VLAN dan konektivitas menuju jaringan inti.

\section{Konsep Implementasi Jaringan Wi-Fi}

\subsection{Service Set Identifier (SSID)}

Subbab ini membahas konsep SSID sebagai identitas jaringan Wi-Fi yang digunakan untuk membedakan layanan jaringan nirkabel.

\subsection{Virtual Local Area Network (VLAN)}

Subbab ini menjelaskan konsep VLAN sebagai mekanisme segmentasi jaringan untuk meningkatkan keamanan, efisiensi, dan pengelolaan jaringan.

\subsection{Trunk dan Access Port}

Subbab ini membahas perbedaan fungsi \textit{access port} dan \textit{trunk port} beserta implementasinya dalam jaringan berbasis VLAN.

\subsection{Dynamic Host Configuration Protocol (DHCP)}

Subbab ini menjelaskan prinsip kerja DHCP dalam proses pemberian alamat IP secara otomatis kepada perangkat klien.

\subsection{Inter-VLAN Routing}

Subbab ini membahas konsep komunikasi antar VLAN melalui mekanisme \textit{routing} sehingga perangkat pada segmen jaringan yang berbeda dapat saling berkomunikasi.

\subsection{Wireless Client Association}

Subbab ini menjelaskan proses \textit{association} antara perangkat klien dan \textit{access point}, mulai dari pemindaian jaringan hingga memperoleh akses layanan.


\section{Pengelolaan dan Konfigurasi WLAN}

\subsection{AP Adoption}

Subbab ini membahas proses registrasi dan integrasi \textit{access point} ke dalam sistem pengelolaan \textit{wireless LAN controller}.

\subsection{Wireless LAN Configuration}

Subbab ini menjelaskan konfigurasi layanan WLAN, termasuk pembuatan SSID, pemetaan VLAN, serta pengaturan parameter jaringan.

\subsection{Radio Configuration}

Subbab ini membahas konfigurasi parameter radio, seperti kanal, lebar pita frekuensi, dan daya pancar untuk mengoptimalkan performa jaringan Wi-Fi.

\subsection{Authentication}

Subbab ini menjelaskan mekanisme autentikasi pengguna pada jaringan Wi-Fi, baik menggunakan \textit{pre-shared key} maupun metode autentikasi lainnya.

\subsection{Firmware Management}

Subbab ini membahas pengelolaan versi \textit{firmware} perangkat jaringan untuk memastikan kompatibilitas, keamanan, dan stabilitas sistem.


\section{Pengujian Jaringan Wi-Fi}

\subsection{Connectivity Testing}

Subbab ini menjelaskan metode pengujian konektivitas untuk memastikan setiap perangkat dapat saling berkomunikasi sesuai dengan rancangan jaringan.

\subsection{Functional Testing}

Subbab ini membahas pengujian terhadap fungsi-fungsi utama jaringan guna memastikan seluruh layanan berjalan sesuai kebutuhan.

\subsection{Coverage Testing}

Subbab ini menjelaskan pengujian cakupan sinyal Wi-Fi untuk mengetahui kualitas layanan pada berbagai area implementasi.

\subsection{Throughput Testing}

Subbab ini membahas pengujian performa jaringan melalui pengukuran kecepatan transfer data pada jaringan Wi-Fi.

\subsection{Site Survey menggunakan NetSpot}

Subbab ini menjelaskan konsep \textit{wireless site survey} menggunakan perangkat lunak NetSpot untuk mengevaluasi kualitas sinyal, cakupan jaringan, dan distribusi layanan Wi-Fi.


\section{Penelitian Terdahulu}

Bagian ini membahas beberapa penelitian terdahulu yang berkaitan dengan implementasi, konfigurasi, maupun pengujian jaringan Wi-Fi \textit{enterprise}. Hasil kajian digunakan sebagai landasan dalam pelaksanaan implementasi dan pengujian skenario jaringan pada APIE Pilot Project Universitas Brawijaya.